\documentclass{article}
\usepackage[utf8]{inputenc}
\usepackage[T1]{fontenc}
\usepackage{geometry}
\usepackage{graphicx}
\usepackage{float}
\usepackage{hyperref}
\usepackage[english,indonesian]{babel}

\title{Optimasi Pemrosesan Template LaTeX Menggunakan Pendekatan Berbasis AI}
\author{Academic Support Agent}
\date{\today}

\begin{document}
\maketitle

\begin{abstract}

\end{abstract}

\chapter{PENDAHULUAN}


\chapter{TINJAUAN PUSTAKA}


% Original content (commented out):
% \documentclass[10pt,a4paper]{article}
\usepackage[utf8]{inputenc}
\usepackage[T1]{fontenc}
\usepackage[spanish]{babel}
\usepackage{amsmath}
\usepackage{amsfonts}
\usepackage{amssymb}
\usepackage{graphicx}
\usepackage{float}
\usepackage{lipsum}
\usepackage{multicol}
\usepackage{xcolor}
\newcommand{\celda}[1]{
	\begin{minipage}{2cm}
		\vspace{2mm}
		#1
		\vspace{2mm}
	\end{minipage}
}
\definecolor{azul}{rgb}{0.36, 0.54, 0.66}
\usepackage[left=2.00cm, right=2.00cm, top=2.00cm, bottom=2.00cm]{geometry}
\author{Jorge Vicente Niño Mocarro}
\begin{document}
	
	\begin{figure}[H]
		\raggedright
		\includegraphics[scale=0.25]{imagenes/unprg.png} \hfill \includegraphics[scale=0.25]{imagenes/faceac.jpg}
	\end{figure}

	\vspace{6mm}
	
	%ESTE CENTER ES EXCLUSIVO PARA EL TITULO DEL PAPER, AUTOR Y UNIVER.
	\begin{center}
		{\Large \textbf{Análisis Situacional de las MYPES frente a la pandemia COVID-19}}\\
		\vspace{2mm}
		{\large Jorge Vicente Niño Mocarro$^{1}$ , Walther Enrique Llenque Sánchez$^{2}$}\\
		\vspace{6.5mm}
		$^1$\textit{Universidad Nacional Pedro Ruiz Gallo. Facultad de Ciencias Económicas, Administrativas y Contables.}\\
		$^2$\textit{Universidad Nacional Pedro Ruiz Gallo. Facultad de Ciencias Económicas, Administrativas y Contables.}
	\end{center}

	\begin{center}
		\textcolor{azul}{\rule{150mm}{0.8mm}}
	\end{center}

      %ESTE ABSTRACT ES PARA EL RESUMEN PROPIAMENTE DICHO Y PARA LAS PALABRAS CLAVES (KEYWORDS) ,NOTA:el comando \par sirve para iniciar el nuevo parrafo con sangría.
	\begin{abstract}
		\lipsum[4] \par \lipsum[5]
		\vspace{2mm}\par
		\underline{\textbf{Palabras Claves:}} \hspace{2mm} \textit{palabra1, palabra2, palabra3.}
	\end{abstract}
	
	\begin{center}
		\textcolor{azul}{\rule{150mm}{0.8mm}}
	\end{center}
	
	\vspace{5mm}

	\begin{multicols}{2}
		\section{Introducción} \label{Intro}
			\lipsum[2] \par \lipsum[1] \par \lipsum[4]
			
		\section{Antecedentes} \label{antece}
			\lipsum[5] \par \lipsum[7]
	\end{multicols}
	
	
	\begin{table}[H]
		\centering
		\begin{tabular}{cccc}
			\hline
			\celda{Nombre} & \celda{Apellidos} & \celda{Edad} & \celda{País}\\
			\hline
			\celda{Jorge} & \celda{Niño} & \celda{22} & \celda{Perú}\\
			\hline
			\celda{Vicente} & \celda{Mocarro} & \celda{23} & \celda{Colombia}\\
			\hline
		\end{tabular}
		\caption{Esta es mi primera tabla en paper}
		\label{tb: Tabla1}
	\end{table}

	\begin{multicols}{2}
		\lipsum[1]
		\section{Consecuencias Mundiales a causa del COVID-19} \label{conse}
			\lipsum
			
			\begin{figure}[H]
				\centering
				\includegraphics[scale=0.4]{imagenes/faceac.jpg}
				\caption{Logo de FACEAC}
				\label{fig: Figura1}
			\end{figure}
		
		\lipsum[1]
		
		\section{Mecanismos Fiscales y Monetarios} \label{mecani}
			\lipsum[5]
		
	\end{multicols}
\end{document}

\end{document}